% 57North Hacklab articles 
\documentclass{article}
\usepackage{graphicx}
\usepackage{palatino}
%\usepackage{draftwatermark}
\title{Articles of 57 North Hacklab Ltd}
\date{22nd Jan 2014}
% Following method of clause numbering is outlined in comp.text.tex
% posting of 16th Feb 2009
% 
% Define a macro to allow the hanging of text from a marker
\makeatletter
\newcommand*{\hangfrom}[1]{%
  \setbox\@tempboxa\hbox{{#1}}%
  \hangindent \wd\@tempboxa
  \noindent\box\@tempboxa}
\makeatother
% Define clause numbering counters
\newcounter{clause}
\renewcommand\theclause{\arabic{clause}}
\newcounter{subclause}[clause]
\renewcommand\thesubclause{\arabic{clause}.\arabic{subclause}}
\newcounter{subsubclause}[subclause]
\renewcommand\thesubsubclause{\arabic{clause}.\arabic{subclause}.\arabic{subsubclause}}
\newcommand\clause[1]{%
  \refstepcounter{clause}
  \par\hangfrom{{\csname theclause\endcsname}\quad }{#1\par}}
\newcommand\subclause[1]{%
  \refstepcounter{subclause}
  \par\hangfrom{\quad{\csname thesubclause\endcsname}\quad }{#1\par}}
\newcommand\subsubclause[1]{%
  \refstepcounter{subsubclause}
  \par\hangfrom{\quad\quad{\csname thesubsubclause\endcsname}\quad }{#1\par}}
% Put a gap between the clauses
\addtolength{\parskip}{\baselineskip}
% Don't number the headings
\setcounter{secnumdepth}{0}
% Shortcut to save some typing
\newcommand{\charityact}{Charities and Trustee Investment (Scotland) Act 2005}
\newcommand{\companyact}{Companies Act 2006}
% End of preamble
\begin{document}
\begin{titlepage}
  \vspace*{\fill}
  \begin{center}
    \includegraphics[width=0.35\textwidth]{logo}\\[1cm]
    \textsc{\LARGE 57 North}\\[1.5cm]
    \textsc{\Large Articles of Association}\\[0.5cm]
    {\large 25th January 2014}
  \end{center}
  \vspace*{\fill}
\end{titlepage}
\section{General}
\subsection{Name}
\clause{The name of the organisation shall be ``57North Hacklab Ltd'', 
hereafter referred to as ``the Hacklab''.}
\subsection{Location}
\clause{The principal office of the Hacklab shall be in Scotland (and
must remain in Scotland).}
\clause{The Hacklab shall hold its regular activities in Aberdeen.}

\section{Interpretation}

\clause{In the articles:}
  \begin{description}
  \item[Address] means a postal address or, for the purposes of
    electronic communication, an e-mail address in each case
    registered with the Hacklab
  \item[Articles] means these articles of association
  \item[Companies Acts] means the Companies Acts (as defined in section 2 of the
  Companies Act 2006), in so far as they apply to the company;
  \item[Director] means a director of the company, and includes any person 
  occupying the position of director, by whatever name called;
  \item[Document] includes, unless otherwise specified, any document sent or 
  supplied in electronic form;
  \item[Electronic Form] has the meaning given in section 1168 of the Companies 
  Act 2006;
  \item[Member] has the meaning given in section 112 of the Companies Act 2006;
  \item[Proxy Notice] has the meaning given in article \ref{clause:proxy};
  \end{description}
Unless the context otherwise requires, other words or expressions contained in 
these articles bear the same meaning as in the Companies Act 2006 as in force on
the date when these articles become binding on the company.

\subsection{Purposes}
\clause{The Hacklab's purposes are:}
\begin{itemize}
\item to promote and and encourage technical, scientific and artistic
  skills and innovation through individual projects, collaboration and
  education;
\item to provide open events to allow the wider community to meet and
  socialise together;
\item to promote and support the use and development of free and open
  technologies, standards, ideas, hardware and software for the
  benefit of all;
\item to provide a shared physical workspace, tools, storage and other
  resources for its members;
\item to work with other bodies with similar or complementary objectives.
\end{itemize}
\subsection{Powers}

\clause{The Hacklab has the power to do anything which is calculated to
further its purposes or is conducive or incidental to doing so.}

\clause{No part of the income or property of the Hacklab may be paid
or transferred (directly or indirectly) to the members - either in the
course of the Hacklab's existence or on dissolution - except where
this is done in direct furtherance of the Hacklab's purposes.}


\subsection{Liability of Members}
\clause{\label{clause:zero_liability}The liability of the members is limited to
a sum not exceeding \pounds1,
being the amount that each member undertakes to contribute to the
assets of the Hacklab in the event of its being wound up while he/she is a 
member or within one year after he/she ceases
  to be a member, for:}
  \begin{enumerate}
  \item
    payment of the company's debts and liabilities incurred before he/she ceases
	to be a member;
  \item
    payment of the costs, charges and expenses of winding up; and
  \item
    adjustment of the rights of the contributories among themselves.
  \end{enumerate}

\subsection{Use of Resources}

\clause{The Hacklab shall make no claim and take no responsibility for
the projects created by users of the Hacklab resources.}

\clause{Use of Hacklab facilities and equipment shall be at the user's
own risk.}

\clause{The Hacklab shall not be held responsible nor liable for
any actions or behaviour of individuals or groups, whether
members or guests.}

\subsection{General Structure}
\clause{The structure of the Hacklab consists of}

  \subclause{the MEMBERS - who have the right to attend members'
  meetings (including any annual general meeting) and have important
  powers under the articles; in particular, the members appoint
  people to serve on the board and take decisions on changes to the
  articles itself;}

  \subclause{the BOARD - who hold regular meetings, and generally
  control the activities of the Hacklab; for example; the board
  is responsible for monitoring and controlling the financial position
  of the Hacklab.}

\clause{The people serving on the board are referred to in these
articles as Directors.}
\section{MEMBERS}

\subsection{General}

\clause{A member shall be permitted to attend any physical workspaces
at any available time, to make use of all Hacklab resources and to accompany
guests.}

\subsection{Qualifications for membership}

\clause{Membership is open to any individual aged 18 or over.}

\subsection{Application for membership}

\clause{Any person who wishes to become a member must submit an 
application for membership, along with a remittance to meet the first
month's subscription; the application will be considered by the
board.}

\clause{The board may, at its discretion, refuse to admit any person
to membership. If membership is refused, the first month's
subscription will be returned.}
 
\clause{The board must notify each applicant promptly (in writing or
by e-mail) of its decision on whether or not to admit him/her to
membership.}

\subsection{Membership subscription}

\clause{Members shall be required to pay a monthly membership
subscription.}

\clause{The monthly membership subscriptions shall be payable on or
before the final day of the month.}

\clause{The members may vary the amount of the monthly membership fee
by way of a resolution to that effect passed at a General Meeting.}

\clause{If the membership subscription payable by any member remains
  outstanding more than 2 months after the date on which it fell due -
  and providing he/she has been given at least one reminder (either in
  writing or by e-mail) - the board may, by resolution to that effect,
  expel him/her from membership.}

\clause{A person who ceases (for whatever reason) to be a member shall
not be entitled to any refund of the membership subscription.}

\subsection{Register of members}

\clause{The board must keep a register of members, setting out:}

\subclause{for each current member:}

\subsubclause{his/her full name and address(postal and email); and}

\subsubclause{the date on which he/she was registered as a member of
the Hacklab.}

\subclause{for each former member - for at least 6 years from the date
on he/she ceased to be a member:}

\subsubclause{his/her name; and}

\subsubclause{the date on which he/she ceased to me a member.}

\clause{The board must ensure that the register of members is updated
within 28 days of any change:}

\subclause{which arises from a resolution of the board or a resolution
passed by the members of the Hacklab; or}

\subclause {which is notified to the Hacklab.}

\clause {Each member is responsible for supplying his/her contact details
to the board to be added to the register of members. Each member is
responsible for ensuring her/his contact details are current by informing
the board of changes when required.}

\subsection{Withdrawal from membership}

\clause{Any person who wants to withdraw from membership must give
  written notice of withdrawal to the Hacklab, signed by him or her,
  or as an e-mail to one or more of the directors; he or she will
  cease to be a member as from the time when the notice is received by
  the Hacklab.}

\clause{Upon ceasing to be a member, all keys, access tokens and
Hacklab property must be returned to the Hacklab.}

\subsection{Transfer of membership}

\clause{Membership of the Hacklab shall not be
> transferrable.}

\subsection{Expulsion from membership}

\clause{\label{clause:expel}Any person may be expelled from membership by way
a resolution passed by not less than two thirds of those present and
voting at a general meeting, providing the following procedures have
been observed:}

\subclause{at least 21 days notice of the intention to propose the
resolution must be given to the member concerned, specifying the
grounds for the proposed expulsion;}

\subclause{the member concerned will be entitled to be heard on the
resolution at the general meeting at which the resolution is
proposed.}

\clause{The board may temporarily suspend any member from access to
  Hacklab resources, if the behaviour of that members is deemed by the
  board to be unacceptable to the Hacklab. Suspension will continue
  until either the board reverses the suspension, or an expulsion
  resolution is heard.}

\subsection{Death of a Member}

\clause{Membership ceases upon the death of the member.}

\section {DECISION MAKING BY THE MEMBERS}

\subsection{General Meetings}

\clause{The board must arrange a meeting of members (an annual general
meeting or ``AGM'') in the month of April each calendar year.}

\clause{The business of each AGM must include:}

\subclause{a report by the chair on the activities of the Hacklab;}

\subclause{consideration of the annual accounts of the Hacklab; and}

\subclause{the election/re-election of directors, as referred
to in clauses \ref{clause:boardelection} to \ref{clause:directorelection}.}

\clause{The board may arrange an extraordinary general meeting at any time.}

\subsection{Power to request the board to arrange a extraordinary
  general meeting (EGM)}

\clause{\label{clause:egm}The board must arrange an extraordinary
general meeting (EGM) if they are requested to do so by a notice
(which may take the form of two or more documents in the same terms,
each signed by one or more members) by members who amount to
20\% or more of the total membership of the Hacklab at the
time, or 5 members which ever is greater providing:}

\subclause{the notice states the purposes for which the meeting is to
be held; and}

\subclause{those purposes are not inconsistent with the terms of these
articles, the \companyact\ or any other statutory provision.}

\clause{If the board receives a notice under clause \ref{clause:egm},
the date for the meeting which they arrange in accordance with the
notice must not be later than 28 days from the date on which they
received the notice.}

\subsection{Notice of general meetings}

\clause{\label{clause:meeting_notice}At least 14 clear days notice must
be given of any AGM or any EGM.}

\clause{The reference to ``clear days'' in clause
\ref{clause:meeting_notice} shall be taken to mean that, in
calculating the days of notice,}

\subclause{the day after the notices are posted (or sent by e-mail)
should be excluded; and}

\subclause{the day of the meeting itself should also be excluded.}

\clause{The notice calling the meeting must specify in general terms
what business is to be dealt with at the meeting; and}

\subclause{in the case of a resolution to alter the articles, must
set out the exact terms of the proposed alteration(s); or}

\subclause{in the case of any other resolution falling within clause
\ref{clause:supermajority} (requirement for a two-thirds majority) must set out
the exact terms of the resolution.}

\clause{Notice of every AGM and EGM must be given to all members and directors
of the Hacklab, but the accidental omission to give notice to one or more members
will not invalidate the proceedings at the meeting.}

\clause{Any notice which requires to be given to a member under these
articles must be sent by e-mail to the member, at the e-mail
address kept on the register of members.}

\subsection{Procedure at AGMs and EGMs}

\clause{No valid decisions can be taken at any AGM or EGM unless a
quorum is present.}

\clause{The quorum for a general meeting is 5 members or one quarter the
membership at the time of the meeting, whichever is greater.}

\clause{If quorum is not present within 15 minutes after the time at
which a general meeting was due to start - or if a quorum ceases to be
present during a general meeting - the meeting cannot proceed; and
fresh notices of a new general meeting must be sent out, to deal with the
business (or remaining business) which was intended to be conducted.}

\clause{A director should act as the chairperson of each general
meeting.}

\subsection{Voting at general meetings}

\clause{Every member has one vote, which can be given personally or by proxy.}

\clause{All decisions at general meetings will be made by majority
vote - with the exception of the types of resolution listed in clause
\ref{clause:supermajority}.}

\clause{\label{clause:supermajority}The following resolutions will be valid only
if passed by not less that two-thirds of those voting on the resolution at an 
AGM or EGM (or if passed by way of a written resolution under clause
\ref{clause:writtenresolution}):}

\subclause{a resolution amending the articles;}

\subclause{a resolution expelling a person from membership under
clause \ref{clause:expel};}

\subclause{a resolution directing the board to take any particular
step (or directing the board not to take any particular step);}

\subclause{a resolution for the winding up or dissolution of the
Hacklab.}

\clause{If there is an equal number of votes for and against any
resolution, the chairperson of the meeting will be entitled to a
second (casting) vote.}

\clause{A resolution put to the vote at a general meeting will be
decided on a show of hands - unless the chair (or at least two other
members present at the meeting) ask for a secret ballot.}

\clause{The chairperson will decide how any secret ballot is to be
conducted, and he/she will declare the result of the ballot at the
meeting.}

\clause{\label{clause:proxy}Proxies may only validly be appointed by a notice in
writing (a “proxy notice”)
which}
\subclause{states the name and address of the member appointing the proxy;}
\subclause{identifies the person appointed to be that member’s proxy and the
general meeting inrelation to which that person is appointed;}
\subclause{is signed by or on behalf of the member appointing the proxy, or is 
authenticated in such manner as the directors may determine; and}
\subclause{is delivered to the company in accordance with the articles and any 
instructions contained in the notice of the general meeting to which they relate.}
\subclause{The company may require proxy notices to be delivered in a particular
 form, and may specify different forms for different purposes.}
\subclause{Proxy notices may specify how the proxy appointed under them is to
vote (or that the proxy is to abstain from voting) on one or more resolutions.}
\subclause{Unless a proxy notice indicates otherwise, it must be treated as—}
\subsubclause{allowing the person appointed under it as a proxy discretion as to
 how to vote on any ancillary or procedural resolutions put to the meeting, and}
\subsubclause{appointing that person as a proxy in relation to any adjournment 
of the general meeting to which it relates as well as the meeting itself.}

\clause{Delivery of proxy notices}

\subclause{A person who is entitled to attend, speak or vote (either on a show
 of hands or on a poll) at a general meeting remains so entitled in respect of
 that meeting or any adjournment of it, even though a valid proxy notice has 
been delivered to the company by or on behalf of that person.}
\subclause{An appointment under a proxy notice may be revoked by delivering to
 the company a notice in writing given by or on behalf of the person by whom 
 or on whose behalf the proxy notice was given.}
\subclause{A notice revoking a proxy appointment only takes effect if it is
delivered before the start of the meeting or adjourned meeting to which it
relates.}
\subclause{If a proxy notice is not executed by the person appointing the proxy,
it must be accompanied by written evidence of the authority of the person who
executed it to execute it on the appointor’s behalf.}

\subsection{Written resolutions by members}

\clause{\label{clause:writtenresolution}A resolution agreed to in
  writing (or by e-mail) by all the members will be as valid as if it
  had been passed at a general meeting; the date of the resolution
  will be taken to be the date on which the last member agreed to it.}

\subsection{Minutes}

\clause{The board must ensure that proper minutes are kept in relation
to all general meetings.}

\clause{Minutes of general meetings must include the names of those
present; and (so far as possible) should be signed by the chairperson
of the meeting.}

\section{Board}

\subsection{Number of directors}

\clause{There shall be 5 directors.}

\subsection{Treasurer}

\clause{One of the directors shall be the treasurer.}

\clause{A person elected to be treasurer will automatically cease to
hold that office:}

\subclause{if he/she ceases to be a director; or}

\subclause{if he/she gives to the Hacklab a notice of resignation
from that office, signed by him/her.}

\subsection{Eligibility}

\clause{A person will not be eligible for
  election to the board unless he/she is a member of the Hacklab.}

\clause{\label{clause:disqualified}A person will not be eligible for election to
the board if he/she is disqualified from being a director under the
\companyact.}

\subsection{Election, retiral, re-election}

\clause{\label{clause:boardelection}At each AGM, the members may elect
any member (unless he/she is debarred from membership under clause
\ref{clause:disqualified}) to be the treasurer or one of the other four
directors.}

\clause{At each AGM, all of the directors must retire from
office - but shall then be
eligible for re-election under clause \ref{clause:boardelection}.}

\clause{A person who has served on the board
for a period of 3 years shall automatically vacate office on expiry of
that 3 year period.}

\clause{\label{clause:directorelection}A director retiring at
  an AGM will be deemed to have been re-elected unless:}

\subclause{he/she advises the board prior to the conclusion of the AGM
that he/she does not wish to be reappointed as a director; or}

\subclause{an election process was held at the AGM and he/she was not
among those elected/re-elected through that process.}

\subsection{Termination of office}

\clause{A director will automatically cease to hold office if:}

\subclause{he/she becomes disqualified from being a director
under the \companyact;}

\subclause{a bankruptcy order is made against him/her;}

\subclause{a composition is made with his/her creditors generally in 
satisfaction of that person’s debts;}

\subclause{he/she becomes incapable for medical reasons of carrying
out his/her duties as a director - but only if that has
continued (or is expected to continue) for a period of more than six
months;}

\subclause{by reason of that person’s mental health, a court makes an order
which wholly or partly prevents him/her from personally exercising any powers
or rights which he/she would otherwise have;}

\subclause{he/she ceases to be a member of the Hacklab;}

\subclause{he/she gives the Hacklab a notice of resignation, signed by
him/her;}

\subclause{he/she is absent (without good reason, in the opinion of
the board) from more than three consecutive meetings of the board -
but only if the board resolves to remove him/her from office.}

\subclause{\label{clause:materialbreach}he/she is removed from office
  by resolution of the board on the grounds that he/she is considered
  to have committed a material breach of the code of conduct for
  directors (as referred to in clause
  \ref{clause:codeofconduct}).}

\subclause{\label{clause:removedbymembers}he/she is removed from
  office by a resolution of the members passed at a general meeting.}

\clause{A resolution under paragraph \ref{clause:materialbreach},
or \ref{clause:removedbymembers} shall
be valid only if:}

\subclause{the director who is the subject of the resolution is
given reasonable prior written notice of the grounds upon which the
resolution for his/her removal is to be proposed;}

\subclause{the director concerned is given the opportunity to
address the meeting at which the resolution is proposed, prior to the
resolution being put to the vote; and}

\subclause{(in case of a resolution under clause
\ref{clause:materialbreach} at least
two thirds (to the nearest round number) of the directors then
in office vote in favour of the resolution.}

\subsection{Register of directors}

\clause{The board must keep a register of directors, setting
out:}

\subclause{for each current directors;}

\subsubclause{his/her full name and address;}

\subsubclause{the date on which he/she was appointed as a directors; and}

\subsubclause{whether he/she is the treasurer.}

\subclause{for each former director - for at least 6 years from
the date on which he/she ceased to be a director:}

\subsubclause{the name of the director;}

\subsubclause{whether he/she was the treasurer; and}

\subsubclause{the date on which he/she ceased to be a director.}

\clause{The board must ensure that the register of directors is
updated within 28 days of any change:}

\subclause{which arises from a resolution of the board or a resolution
passed by members of the Hacklab; or}

\subclause{which is notified to the Hacklab.}


\subsection{Powers of board}

\clause{Except where these articles state otherwise, the
Hacklab (and its assets and operations) will be managed by the
board; and the board may exercise all the powers of the Hacklab.}

\clause{A meeting of the board at which a quorum is present may
exercise all powers exercisable by the board.}

\clause{The members may, by way of a resolution passed in compliance
with clause \ref{clause:supermajority} (requirement for two-thirds
majority), direct the board to take any particular step or direct the
board not to take any particular step; and the board shall give effect
to any such direction accordingly.}

\subsection{Directors - general duties}

\clause{\label{clause:directorduties}Each of the directors has a
  duty, in exercising functions as a director, to act in the
  interests of the Hacklab; and in particular, must:}

\subclause{seek, in good faith, to ensure that the Hacklab acts
  in a manner which is in accordance with its purposes;}

\subclause{act with the care and diligence which it is reasonable to
expect of a person who is managing the affairs of another person;}

\subclause{in circumstances giving rise to the possibility of a
conflict of interest between the Hacklab and any other party:}

\subsubclause{put the interests of the Hacklab before that of the
other party;}

\subsubclause{where any other duty prevents him/her from doing so,
disclose the conflicting interest to the Hacklab and refrain from
participating in any deliberation or decision of the other directors
with regard to the matter in question;}

\subclause{ensure that the Hacklab complies with any direction,
requirement, notice or duty imposed under or by virtue of the
\companyact.}

\clause{In addition to the duties outlined in clause
  \ref{clause:directorduties}, all of the directors must take
  such steps as are reasonably practicable for the purpose of
  ensuring:}

\subclause{that any breach of any of those duties by director
  is corrected by the director concerned and not repeated; and}

\subclause{that any director who has been in serious and persistent
  breach of those duties is removed as a director.}

\clause{Provided he/she has declared his/her interest - and has not
  voted on the question of whether or not the Hacklab should enter
  into the arrangement - a director will not be debarred from
  entering into an arrangment with the Hacklab in which he/she has a
  personal interest and he/she may retain any personal benefit which
 arises from that arrangement.}

\clause{\label{clause:directoremployee}No director may be given 
any remuneration by the Hacklab for carrying out his/her duties as a director.}

\subsection{Code of conduct for directors}

\clause{\label{clause:codeofconduct}Each of the directors shall
  comply with the code of conduct (incorporating detailed rules on
  conflict of interest) prescribed by the board from time to
  time.}

\clause{The code of conduct referred to in clause
  \ref{clause:codeofconduct} shall be supplemental to the conduct of
  director contained in the articles and the duties
  imposed on directors under the \companyact; and all relevant
  provisions of this consitution shall be interpreted and applied in
  accordance with the provisions of the code of conduct in force from
  time to time.}

\section{DECISION-MAKING BY THE DIRECTORS}

\subsection{Notice of board meetings}

\clause{Any director may call a meeting of the board.}

\clause{At least 7 days' notice must be given of each board meeting,
  unless (in the opinion of the person calling the meeting) there is a
  degree of urgency which makes that inappropriate.}

\subsection{Procedure at board meetings}

\clause{\label{clause:boardquorum}No valid decisions can be taken at a
  board meeting unless a quorum is present; the quorum for board
  meetings is 3 directors, present in person.}

\clause{If at any time the number of directors in office falls
  below the number stated as the quorum in clause
  \ref{clause:boardquorum}, the remaining directors will have
  power to call an EGM - but will not be able to take any other valid
  decisions.}

\clause{At each board meeting, the directors present must elect
  (from amongst themselves) the person who will act as chairperson of
  that meeting.}

\clause{Each directors has one vote, which must be given in
  person.}

\clause{All decisions at board meetings will be made by majority
  vote.}

\clause{If there is an equal number of votes for and against any
  resolution, the chairperson of the meeting will be entitled to a
  second (casting) vote.}

\clause{The board may, at its discretion, allow any person to attend
  and speak at a board meeting notwithstanding that he/she is not a
  director - but on the basis that he/she must not participate
  in decision-making.}

\clause{\label{clause:conflictofinterest}A director must not
  vote at a board meeting (or at a meeting of a sub-committee) on any
  resolution which relates to a matter in which he/she has a personal
  interest or duty which conflicts (or may conflict) with the
  interests of the Hacklab; he/she must withdraw from the meeting
  while an item of that nature is being dealt with.}

\clause{For the purposes of clause \ref{clause:conflictofinterest}: 
  a director will be deemed to have a personal
  interest in relation to a particular matter if a body in relation to
  which he/she is an employee, director, member of the management
  committee, officer or elected representative has an interest in that
  matter.}

\subsection{Minutes}

\clause{\label{clause:properminutes}The board must ensure that proper
  minutes are kept in relation to all board meetings and meetings of
  sub-committees.}

\clause{The minutes to be kept under clause \ref{clause:properminutes}
  must include the names of those present; and (so far as possible)
  should be signed by the chairperson of the meeting. They must be
  circulated to all members by email within 30 days of the board meeting.}

\section{Administration}

\subsection{Delegation to sub-committees}

\clause{\label{clause:delegation}The board may delegate any of their
  powers to sub-committees; a sub-committee must include at least one
  director, but other members of a sub-committee need not be
  directors.}

\clause{When delegating powers under clause \ref{clause:delegation},
  the board must set out appropriate conditions (which must include an
  obligation to report regularly to the board).}

\clause{Any delegation of powers under clause \ref{clause:delegation}
  may be revoked or altered by the board at any time.}

\clause{The rules of procedure for each sub-committee, and the
  provisions relating to membership of each sub-committee, shall be
  set by the board.}

\subsection{Operation of accounts}

\clause{Five account signatories will be appointed by the board;
  these must themselves all be directors.}

\clause{\label{clause:signatories}Subject to clause
  \ref{clause:electronicbanking}, the signatures of two out of five 
  signatories will be required in relation to all operations (other
  than the lodging of funds) on the bank and building society accounts
  held by the Hacklab.}

\clause{\label{clause:electronicbanking}Where the Hacklab uses
electronic facilities for the operation of any bank or building
society account, the authorisations required for operations on that
account must be consistent with the approach reflected in clause
\ref{clause:signatories}.}

\subsection{Accounting records and annual accounts}

\clause{The board must ensure that proper accounting records are kept,
  in accordance with all applicable statutory requirements.}

\clause{The board must prepare annual accounts, complying with all
  relevant statutory requirements.}

\section{Miscellaneous}

\subsection{Winding-up}

\clause{Any surplus assets available to the Hacklab immediately
preceding its winding up or dissolution must be used for purposes
which are the same as - or which closely resemble - the purposes of
the Hacklab as set out in these articles.}

\subsection{Alterations to the articles}

\clause{These articles may be altered by resolution of the members
passed at an AGM or EGM (subject to achieving the two-thirds majority
referred to in clause \ref{clause:supermajority}) or by way of a
written resolution of the members.}

\end{document}
